
\begin{abstract}
This survey maps the transition from cascaded ASR→LLM→TTS pipelines to unified Speech Language Models (SLMs) that reason and speak natively, lowering latency, preserving paralinguistics, and enabling duplex behaviors. We organize the modern speech‑AI stack around token interfaces and a representation–latency–bitrate trilemma, synthesizing a taxonomy spanning semantic, acoustic, paralinguistic, mixed, and continuous pathways; quantization mechanisms (k‑means, VQ/RVQ, BSQ); and rate/sequence controls (variable frame‑rates, dedup/acoustic BPE). Architecturally, we review adapter bridges that fuse frozen SSL encoders to text LLMs, decoder‑only SLMs with AR/NAR schedules (e.g., AR‑on‑first RVQ stream with NAR refinement), and interleaved speech–text modeling scaled by token‑space pretraining, highlighting a robust 12.5 Hz operating point and ultra‑low‑rate viability at 6.25 Hz with text‑aware diffusion. We systematize streaming and duplex design—chunked policies, wait‑k/CIF gating, barge‑in/IRQ control—and connect them to practical corpora (e.g., J‑CHAT) and bandwidth ladders. Rate‑aware evaluation is consolidated across datasets (LibriSpeech, Common Voice, CoVoST2, LibriTTS, VoxCeleb) and suites (DASB, STAB, Dynamic‑SUPERB, VoxEval), with metrics for semantics (WER/CER, BLEU), perception (MOS/UTMOS, FAD), unit discriminability (ABX), identity/style, and streaming latency (RTF, AL/LAAL). Contributions include: (i) a unified token taxonomy with explicit rate/bitrate attributes; (ii) a comparative design analysis linking quantization, codebook usage, and rate control to intelligibility, naturalness, and style; and (iii) standardized, bitrate‑controlled training/evaluation practices for reproducible cross‑paper comparison. We close with open challenges—unified vs multi‑stream tokens, long‑context conversational memory, multilingual scaling, alignment without forgetting, and a rate–capacity theory—and provide recommendations for duplex‑ready, safe, and efficient deployment.
\end{abstract}
